\chapter{The Special Period}
\label{ch:collapse}

\section{The stop}

What happened to Cuba between 1990 and 1993 has no close parallel among
countries not at war.

The Soviet bloc had taken roughly 85 per cent of Cuban trade. Comecon dissolved
in June 1991; the Soviet Union followed in December. The preferential sugar
protocols ended, the concessional oil ended, the spare parts ended, the credit
lines ended, and the machinery of an entire economy --- built over eighteen
years to Soviet standards, with Soviet inputs, for a Soviet market ---
discovered simultaneously that none of its three sides existed any more.

Cuban imports fell from about \$8.1 billion in 1989 to about \$2.2 billion in
1993, a fall of roughly three quarters in four years.\unv{} Oil imports fell
from around 13 million tonnes to around 6.\unv{} The official Cuban figure for the fall
in output is 34.8 per cent between 1989 and 1993, probably the largest
four-year decline in twentieth-century Cuba. It is also the \emph{lowest}
figure in circulation: Cuban media reported 48 per cent at the time, and
independent reconstructions have gone as high as 50 per cent
\citep{mesalago2005}. The spread turns
on the exchange-rate question of the note on sources rather than on any dispute
about what physically happened, and this book uses the official figure because
it is the most conservative one available and the argument does not need a
larger number. For comparison, the United States lost about 27 per cent of
output in the Great Depression, over four years, with its trading partners
still in place.

Fidel Castro named it in 1990: the Special Period in Time of Peace, a phrase
lifted from the civil-defence contingency plans drawn up in the 1980s for the
event of a war and a total blockade. The name was accurate. The plans had been
written for a siege, and a siege is what arrived, except that nobody was
besieging.

\section{What that was like}

The abstraction is useless without the texture, and the texture is what Cubans
who lived through it still measure everything against.

Electricity went first, because the thermal plants ran on imported fuel.
Blackouts of twelve to sixteen hours a day were normal by 1993, scheduled by
neighbourhood and then not scheduled at all. Without electricity there is no
water, because Cuban water is pumped; no refrigeration, in a climate where food
spoils in hours; no lifts in the apartment blocks the microbrigades had built;
and no light to read by in a country whose principal remaining asset was that
everyone could read.

Transport went next. Buses stopped for want of fuel and tyres. The state
imported over a million Chinese bicycles and built factories to assemble more,
and Havana --- a city with no bicycle tradition at all --- became a city of
bicycles, with the injuries and the exhaustion that implies in that heat. Later
came the \emph{camellos}, articulated trailers hauled by lorry tractors,
carrying several hundred standing passengers each. Oxen returned to fields that
had been mechanised for twenty years, and the state trained them and the men to
drive them, because there was no diesel.

Then food. The ration book had always been a floor and the state economy a
ceiling; now the ceiling collapsed onto the floor and the floor did not hold.
Average daily energy intake fell from roughly 2,900 kilocalories in 1989 to
somewhere around 1,860 in 1993, and protein intake fell by close to half.\unv{}
Cuban adults lost, on average, between five and ten kilograms. There was no
famine in the technical sense --- no mass starvation deaths, no localised
collapse --- but there was a country-wide, sustained, involuntary reduction in
food intake to below what an adult needs.

The clinical proof arrived in 1992. An epidemic of optic and peripheral
neuropathy spread across the island, blurring and in some cases destroying
vision. It ran from 1991 to 1994 and peaked over 1992--93, and the Cuban health
ministry's own count was 50,862 cases in a population of 10.8 million
\citep{mmwr1994}. The
cause was nutritional: B-vitamin deficiency, aggravated by tobacco and rough
alcohol, in a population that had been eating too little for two years. It was
halted by handing out vitamin tablets to the entire population, which is
simultaneously an indictment of the situation and a demonstration of what a
health system that reaches every block can do in a month. There were no deaths
from it, and fewer than one in a thousand of those affected were left with
lasting damage --- which is the same story as the infant mortality figures
below, and should be read the same way.

Everything else followed from those four. Soap and detergent disappeared, and
with them a measurable rise in diarrhoeal and skin disease. Tuberculosis
incidence roughly tripled from a very low base. Prostitution reappeared, aimed
at the tourists who were the only source of hard currency, and with it a word,
\emph{jineterismo}, that had not been needed for thirty years. The sugar
industry --- the thing the whole economy had been built around --- could not get
fertiliser, fuel or spares, and the harvest fell from 8.1 million tonnes in
1989--90 to 4.3 million by 1993 and 3.3 million by 1994--95. It never
recovered, in any year, ever again.

\section{What did not collapse}

Here the book has to be careful, because this is the single strongest piece of
evidence in favour of the Cuban model and it is routinely omitted by people who
find it inconvenient.

Infant mortality did not rise. It fell. In 1989 it was about 11.1 per thousand
live births; through the worst years of the Special Period it continued to
decline, and by 1994 it was under 10.\unv{} Life expectancy did not fall.
Immunisation coverage did not fall. No school closed. The universities stayed
open, on bicycles and candlelight. There was no epidemic of the classic
mortality signatures of economic collapse --- no surge in infant deaths, no
collapse of male life expectancy of the kind Russia recorded in exactly the
same years from exactly the same shock.

This is not a small thing and it should not be explained away. Russia in
1990--94 lost a comparable share of output and lost roughly six years of male
life expectancy with it. Cuba lost more output and lost none. The difference is
that Cuba's health and education systems were not financed out of the collapsing
sector; they were a prior claim on whatever remained, defended politically, and
staffed by people who kept coming to work in a country with no transport. When
this book argues in Chapter~\ref{ch:welfare} that the welfare state is the one
thing worth keeping, this is the evidence: a universal system with deep
territorial reach is a shock absorber of a kind that nothing else in a poor
country's toolkit can substitute for.

Two qualifications, both real. The systems degraded even where the outcome
indicators held --- hospitals without anaesthetics, without X-ray film, without
functioning sterilisers, and doctors who improvised or did without. And there
was one genuinely perverse benefit, documented afterwards and worth stating
because it shows how strange the episode was: the forced calorie restriction and
the forced physical activity produced a measurable fall in diabetes and
cardiovascular mortality across the 1990s, which reversed when weight was
regained after 1995 \citep{franco2013}.\unv{} A country lost weight because it was hungry and its
heart-disease deaths fell. That is not a policy recommendation. It is a reminder
that this period does not fit any simple story.

\section{Option Zero, and what came out of it}

The contingency planning went further than what was implemented. Option Zero was
the plan for the total loss of imported fuel: no electricity generation to
speak of, transport by animal, food produced and distributed at neighbourhood
level through communal kitchens, the country reorganised into self-supplying
cells. It was prepared and it was not needed, and that it was prepared says
something about how the leadership assessed 1992.

What was implemented, and what outlasted the crisis, was the agricultural
improvisation. Without fertiliser, pesticide or fuel, Cuban agriculture was
forced into low-input methods: biological pest control, composting, animal
traction, and above all the \emph{organop\'onicos}, raised-bed vegetable
gardens on vacant urban lots, which by the 2000s supplied a substantial share
of the fresh vegetables consumed in Havana. This became internationally famous
and is frequently over-claimed. What it demonstrably did was put vegetables into
cities. What it did not do was feed the country: Cuba's food import dependence
went up across the whole period, not down, and the case for low-input
agriculture has to be made on its merits in Chapter~\ref{ch:land} rather than
on a 1990s legend.

\section{August 1994}

By the summer of 1994 the currency had gone. The peso, official at one to the
dollar, traded on the street at 120 to 150. A state wage bought almost nothing.
People were leaving on anything that floated, and the state's practice of
intercepting boats had produced a series of violent incidents.

On 13 July 1994 a group of about seventy people hijacked an old tugboat, the
\emph{13 de Marzo}, and headed for Florida. It was pursued by state-operated
vessels which rammed it and directed high-pressure water hoses at the people on
deck, including at children being held up. The tug sank seven miles offshore.
Thirty-seven people died, ten of them children.\unv{} The government's account
was that the sinking was accidental and the fault of the hijackers. No
independent investigation was permitted. The Inter-American Commission on Human
Rights found the state responsible \citep{iachr1996}. Whatever one concludes about the Special
Period generally, this happened, and it happened to people trying to leave a
country that would not let them.

On 5 August 1994, after rumours that a boat would be allowed to depart,
thousands of people gathered on the Havana Malec\'on, threw stones, broke
windows, and shouted against the government. It was the first mass anti-
government demonstration in thirty-five years. It was put down within hours,
partly by police and partly by mobilised construction workers with sticks, and
Castro himself came to the seafront during it. The \emph{maleconazo} lasted an
afternoon and changed the government's calculation permanently.

Within days Castro announced that anyone who wanted to leave by sea could go
unhindered. Over the following five weeks something like 30,000 to 35,000 people
put to sea on rafts and inner tubes.\unv{} An unknown number drowned. The United
States, unable to absorb them, changed a thirty-year-old policy: rafters
intercepted at sea were taken to Guant\'anamo rather than to Florida, where
some 30,000 were held in camps for months. The negotiations that followed
produced the migration accords of September 1994 and May 1995, under which the
United States undertook to admit a minimum of 20,000 Cubans a year through
legal channels and to return those interdicted at sea, while those who reached
land could stay --- the policy that became known as wet-foot, dry-foot, and that
would shape Cuban emigration for the next twenty-two years.

The migration valve is the mechanism this book will keep returning to. Twice
before --- Camarioca in 1965, Mariel in 1980 --- and now a third time, internal
pressure was released by opening the door, and each time it worked. It is one of
the two instruments that have kept the Cuban political system stable through
crises that would otherwise have broken it, the other being the welfare floor.
Chapter~\ref{ch:crisis} will show the same valve opening a fourth time, on a
scale that finally threatened the thing it was protecting.

\begin{ledger}{the Special Period, 1990--1994}
\emph{Achieved.} A loss of over a third of national output absorbed without
famine, without epidemic mortality, without the collapse of the state, and
without the fall in life expectancy that the identical shock produced in
Russia. Infant mortality continued to \emph{fall} throughout. Every school
stayed open. The health system degraded severely and held its outcomes, which
is the strongest evidence in this book that a universal system with
territorial reach is worth the money it costs. Urban agriculture, born of
desperation, became a durable and genuinely useful institution.

\emph{Cost.} A country of readers reduced to roughly 1,860 kilocalories a day,
losing five to ten kilograms a head, with fifty thousand cases of nutritional
blindness as the clinical receipt. The sugar industry destroyed and never
recovered. A decade of consumption gone. Prostitution and a dollar economy
restored as facts of Cuban life. Thirty-seven people, ten of them children,
drowned off Havana in an incident the state has never independently
investigated. And the moment, on an August afternoon on the Malec\'on, when the
social contract of 1959 stopped being believed by the people it had been
written for.
\end{ledger}
